\documentclass[10pt,a4paper]{altacv}
\usepackage{cv-style}
\begin{document}
    \name{Angela Cortecchia}
    \tagline{Ph.D. student}
    \personalinfo{%
        % You can add your own with \printinfo{symbol}{detail}
        \email{angela.cortecchia@unibo.it}
        \email{angela.cortecchia@hotmail.com}
    %   \phone{+49-173-6895039} 
        \\
        \github{github.com/angelacorte}
        \linkedin{linkedin.com/in/angela-cortecchia}
        \\
        \location{Cesena, Italy}

        %\orcid{orcid.org/0000-0000-0000-0000} % Obviously making this up too. If you want to use this field (and also other academicons symbols), add "academicons" option to \documentclass{altacv}
    }

    %% Make the header extend all the way to the right, if you want. Extend the right margin by 8cm (=6.8cm marginparwidth + 1.2cm marginparsep)
    \begin{adjustwidth}{}{-8cm}
        \makecvheader
    \end{adjustwidth}

\begin{adjustwidth}{}{-9cm}
    \cvsection{Research Profile}
    Ph.D. student in Computer Science and Engineering working on collective robotics through
    macro-programming. My research focuses on aggregate computing, distributed coordination,
    and software abstractions for heterogeneous multi-robot systems. I design and evaluate
    research software, primarily with Kotlin Multiplatform, for reproducible experiments on
    self-organizing collective behaviors.
\end{adjustwidth}

    %%% EDUCATION
    \begin{adjustwidth}{}{-9cm}
    \cvsection{Education}
    \small

    \textcolor{accent}{\textbf{Ph.D. in Computer Science and Engineering}}, University of Bologna
    \hfill \textit{Nov 2024--ongoing, Cesena, Italy}\\
    \emph{Advances in Collective Robotics through Macro-programming}\\
    \textbf{Supervisor:} Danilo Pianini \quad
    \textbf{Co-supervisor:} Mirko Viroli

    \divider

    \textcolor{accent}{\textbf{M.Sc. in Engineering and Computer Science}}, University of Bologna
    \hfill \textit{2021--2024, Cesena, Italy}\\
    Final grade: \textbf{109/110}\\
    \emph{A Kotlin Multiplatform Implementation of Aggregate Computing Based on XC}\\
    \textbf{Supervisor:} Danilo Pianini \quad
    \textbf{Co-supervisor:} Nicolas Farabegoli

    \divider

    \textcolor{accent}{\textbf{B.Sc. in Engineering and Computer Science}}, University of Bologna
    \hfill \textit{2017--2021, Cesena, Italy}\\
    \emph{HCI Methodologies for Developing an Online Cultural Events App}\\
    \textbf{Supervisor:} Silvia Mirri

    \divider

    \textcolor{accent}{\textbf{Undergraduate Internship}}, University of Bologna
    \hfill \textit{Oct--Dec 2020, Cesena, Italy}\\
    Development of an online cultural-events application.\\
    \textbf{Supervisor:} Silvia Mirri

    \divider

    \textcolor{accent}{\textbf{Technical Diploma in Business Administration and Foreign Languages}},
    ITS ``A. Oriani''
    \hfill \textit{2017, Faenza, Italy}

    \normalsize
    \end{adjustwidth}

    %%% PUBLICATIONS
    \begin{adjustwidth}{}{-9cm}
    \cvsection{Publications}

    \cvgroupheading{Journal Articles}
    \nocite{*}
    \printbibliography[
        heading=none,
        type=article,
        check=pub
    ]
    \printbibliography[
        heading=none,
        type=article,
        check=notpub
    ]

    \divider

    \cvgroupheading{Conference Papers}
    \printbibliography[
        heading=none,
        type=inproceedings,
        notkeyword=workshop,
        check=pub
    ]
    \printbibliography[
        heading=none,
        type=inproceedings,
        notkeyword=workshop,
        check=notpub
    ]

    \divider

    \cvgroupheading{Workshop Papers}
    \printbibliography[
        heading=none,
        type=inproceedings,
        keyword=workshop,
        check=pub
    ]
    \printbibliography[
        heading=none,
        type=inproceedings,
        keyword=workshop,
        check=notpub
    ]

    \divider
    \end{adjustwidth}

%%% RESEARCH EXPERIENCE

    \begin{adjustwidth}{}{-9cm}
        \cvsection{Research Experience}
        \small

        \textcolor{accent}{\textbf{Research Fellow}}, University of Bologna -- DISI / Consortium GARR
        \hfill \textit{15 Feb--31 Oct 2024, Cesena, Italy}\\
        \emph{A Unifying Approach to Programming Heterogeneous Devices
        in the Edge-Cloud Continuum}\\
        \textbf{Supervisor:} Danilo Pianini

        \normalsize
    \end{adjustwidth}


%%% SCIENTIFIC ACTIVITIES

    \begin{adjustwidth}{}{-10cm}
    \cvsection{Scientific Activities}
    \small

    \cvgroupheading{Awards and Grants}

    \textcolor{accent}{\textbf{Best Student Paper Award}}, ACSOS 2025
    \hfill \textit{29 Sep 2025, Tokyo, Japan}\\
    \emph{A Field-based Approach for Runtime Replanning in Swarm Robotics Missions}
    \par\smallskip

    \textcolor{accent}{\textbf{Participation Grant}}, ACSOS 2024
    \hfill \textit{16 Sep 2024, Aarhus, Denmark}

    \divider

    \cvgroupheading{Academic Service}

    \textcolor{accent}{\textbf{Publicity Chair}}, ACSOS 2026
    \hfill \textit{7--11 Sep 2026, Cesena, Italy}\\
    \textcolor{accent}{\textbf{Artifact Evaluation Committee Member}}, FormaliSE 2026
    \hfill \textit{12--13 Apr 2026, Rio de Janeiro, Brazil}\\
    \textcolor{accent}{\textbf{Local Chair Support}}, DS-RT 2024
    \hfill \textit{7--9 Oct 2024, Urbino, Italy}

    \divider

    \cvgroupheading{Conference Participation and Presentations}

    \textcolor{accent}{\textbf{Main Track Presenter}}, DCOSS-IoT 2026
    \hfill \textit{22--24 Jun 2026, Reykjavík, Iceland}\\
    \textcolor{accent}{\textbf{Dissemination Track Presenter}}, WOA 2026
    \hfill \textit{15--17 Jun 2026, Salerno, Italy}\\
    \textcolor{accent}{\textbf{Attended}}, COORDINATION at DisCoTec 2026
    \hfill \textit{8--12 Jun 2026, Urbino, Italy}\\
    \textcolor{accent}{\textbf{PhD Symposium Presenter}}, DS-RT 2024
    \hfill \textit{7--9 Oct 2024, Urbino, Italy}\\
    \textcolor{accent}{\textbf{Main Track and PhD Symposium Presenter}}, ACSOS 2024
    \hfill \textit{16--20 Sep 2024, Aarhus, Denmark}

    \divider

    \cvgroupheading{Challenges and Doctoral Schools}

    \textcolor{accent}{\textbf{Participant}}, SPACERAISE 2025
    \hfill \textit{12--16 May 2025, L'Aquila, Italy}\\
    \textcolor{accent}{\textbf{Participant}}, BISS 2025
    \hfill \textit{19--23 May 2025, Bertinoro, Italy}\\
    \textcolor{accent}{\textbf{Participant}}, SIESTA 2025
    \hfill \textit{27--29 Aug 2025, Lugano, Switzerland}\\
    \textcolor{accent}{\textbf{Flash Talk Challenge}}, RFTC 2024
    \hfill \textit{11 Oct 2024, Bologna, Italy}

    \normalsize
    \end{adjustwidth}

    %%% TEACHING EXPERIENCE
    \begin{adjustwidth}{}{-9cm}
        \cvsection{Teaching Experience}
        \small

        \textcolor{accent}{\textbf{Teaching Tutor}}, Programmazione ad Oggetti
        \hfill \textit{Sep 2025--ongoing, Cesena, Italy}\\
        Bachelor's Degree in Computer Science and Engineering,
        University of Bologna

        \divider

        \textcolor{accent}{\textbf{Teaching Tutor}}, Architetture degli Elaboratori
        \hfill \textit{Jan--Sep 2025, Cesena, Italy}\\
        Bachelor's Degree in Computer Science and Engineering,
        University of Bologna

        \normalsize
    \end{adjustwidth}

%%% THESIS CO-SUPERVISION

    \begin{adjustwidth}{}{-9cm}
    \cvsection{Thesis Co-Supervision}
    \small

    \setlength{\tabcolsep}{0pt}
    \renewcommand{\arraystretch}{1.18}

    \cvgroupheading{2026}

    \begin{tabularx}{\linewidth}{@{}
            >{\bfseries\raggedright\arraybackslash}p{3.15cm}
            >{\raggedright\arraybackslash}X@{}}

        Master's thesis &
        \href{https://amslaurea.unibo.it/id/eprint/38087/}
        {Lowering the Reality Gap in Aggregate Programs Validation:
        Running Collektive Over Unity}
        \\

        Bachelor's thesis &
        Un linguaggio per domarli tutti: applicazioni full stack
        con librerie native in Kotlin Multiplatform

    \end{tabularx}

    \divider

    \cvgroupheading{2025}

    \begin{tabularx}{\linewidth}{@{}
            >{\bfseries\raggedright\arraybackslash}p{3.15cm}
            >{\raggedright\arraybackslash}X@{}}

        Master's thesis &
        \href{https://amslaurea.unibo.it/id/eprint/34470/}
        {Static Analysis of Aggregate Programs through Compiler Plugins}
        \\

        Bachelor's thesis &
        \href{https://amslaurea.unibo.it/id/eprint/37130/}
        {Da MIT Proto a Collektive: Trasposizione di esempi di
        programmazione aggregata in un DSL in Kotlin}
        \\

        Bachelor's thesis &
        Multi-platform distributed systems with Aggregate Computing in Kotlin: the case of proximity messaging

    \end{tabularx}

    \normalsize
\end{adjustwidth}


%\begin{adjustwidth}{}{-9cm}
%    \cvsection{Previous work experience}
%
%    \cvblock{Horse Rider and Coach}{Delta Team}{2020--2024}{Alfonsine, Italy}
%    Practicing horse riding at an international competitive level for several years, training and competing young horses. \smallbreak
%    %operatore ludico
%    Done the coaching level ``Operatore Ludico'' of the Italian Equestrian Federation (``Federazione Italiana Sport Equestri'' - FISE) coaching course in 2021. \smallbreak
%
%    \divider
%
%\end{adjustwidth}

    %%% TECHNICAL SKILLS
    \begin{adjustwidth}{}{-9cm}
    \cvsection{Technical Skills}
    \small

    \cvgroupheading{Research and Scientific Computing}

    \textbf{Aggregate and distributed systems:}
    aggregate computing, field calculus, macro-programming,
    collective robotics, concurrent and distributed programming\\

    \textbf{Research software:}
    Collektive, ScaFi, Alchemist, Kotlin Multiplatform, Akka Actors\\

    \textbf{Scientific computing and data analysis:}
    Python, NumPy, pandas, Matplotlib, Seaborn, SQL, MongoDB

    \divider

    \cvgroupheading{Software Development}

    \textbf{Programming languages:}
    Kotlin, Scala, Python, TypeScript, Java, Rust, C, Bash, Prolog\\

    \textbf{Software engineering:}
    software architecture, API design, middleware, microservices,
    automated testing, simulation pipelines\\

    \textbf{Testing and build tools:}
    JUnit, Kotest, Gradle

    \divider

    \cvgroupheading{Development Infrastructure and Applications}

    \textbf{Reproducibility and CI/CD:}
    Git, Docker, GitHub Actions, LaTeX\\

    \textbf{Web and application technologies:}
    Node.js, Express, Android, Angular, React, Svelte, Vue\\

    \textbf{Markup and data formats:}
    HTML, CSS, Markdown, YAML, JSON, XML

    \normalsize
    \end{adjustwidth}

    %%% RESEARCH INTERESTS
\begin{adjustwidth}{}{-9cm}
    \cvsection{Research Interests}

    \cvtag{Aggregate Computing}
    \cvtag{Collective Robotics}
    \cvtag{Self-organizing Systems}
    \cvtag{Collective Intelligence}
    \cvtag{Distributed Systems}
    \cvtag{Macro-programming}
    \cvtag{Reproducible Research}
    \cvtag{Multi-Robot Systems}
    \cvtag{Astronomy}

    \divider
\end{adjustwidth}

\begin{adjustwidth}{}{-9cm}
    \cvsection{Software Projects}

    \cvgroupheading{Research Software}

    \textcolor{accent}{\textbf{Collektive}}\\
    \github{\href{https://github.com/Collektive}{github.com/Collektive}}\\
    \textit{10/2023--ongoing}\\ \smallskip
    Collektive: A Kotlin Multiplatform implementation of Aggregate Computing based on XC. The project was extended for academic purposes for the Master Degree Thesis.\\ \smallskip
%    {\small \notesymbol \hspace{0.5em} Paper published to EUMAS 2023}\\
    \textbf{Language:} Kotlin\\
    \textbf{Keywords:} Domain Specific Language, Aggregate Computing, Field Calculus, eXchange Calculus

   \divider

    \textcolor{accent}{\textbf{CAROL: Coordinated Aggregate Robotics with Online control Lyapunov and barrier functions}}\\
    \github{\href{https://github.com/angelacorte/carol-experiments}{github.com/angelacorte/carol-experiments}}\\
    \textit{01/2026--ongoing}\\ \smallskip
    CAROL is a fully distributed, multi-robot control framework written in Kotlin. It bridges the gap between Aggregate Programming (via the Collektive framework) and Optimization-based Control using Control Lyapunov Functions (CLFs) and Control Barrier Functions (CBFs).\\ \smallskip
%    {\small \notesymbol \hspace{0.5em} Paper published to EUMAS 2023}\\
    \textbf{Language:} Kotlin\\
    \textbf{Keywords:} Aggregate Computing, Control Lyapunov Functions, Control Barrier Functions, Multi-Robot Systems

   \divider

   \textcolor{accent}{\textbf{FieldVMC}}\\
   \github{\href{https://github.com/angelacorte/fieldVMC}{github.com/angelacorte/fieldVMC}}\\
   \textit{02--06/2025}\\ \smallskip
   Experimental artifact for \emph{FieldVMC}, a field-based generalisation of the VMC model for
   self-organizing spatial structures, developed for the related journal article.\\ \smallskip
   {\small \notesymbol \hspace{0.5em} Paper in phase of publishing at the Journal ``Complex And Intelligent Systems''}\\
   \textbf{Language:} Kotlin\\
   \textbf{Keywords:} Aggregate Computing, Self-Organizing, Vascular Morphogenesis Controller

    \divider

    \textcolor{accent}{\textbf{VMC-experiments}}\\
    \github{\href{https://github.com/angelacorte/vmc-experiments}{github.com/angelacorte/vmc-experiments}}\\
    \textit{03--05/2024}\\ \smallskip
    Experimental artifact evaluating a field-based generalisation of the VMC model through
    Aggregate Programming, developed for the related ACSOS 2024 paper.\\ \smallskip
    {\small \notesymbol \hspace{0.5em} Paper published to ACSOS 2024}\\
    \textbf{Language:} Kotlin\\
    \textbf{Keywords:} Aggregate Computing, Self-Organizing, Vascular Morphogenesis Controller

    \divider

    \textcolor{accent}{\textbf{Collektive-examples}}\\
    \github{\href{https://github.com/Collektive/collektive-examples}{github.com/Collektive/collektive-examples}}\\
    \textit{01--03/2024}\\ \smallskip
    Examples of the Collektive project. The examples were made for academic purposes for the Master Degree Thesis and Research Grant.\\ \smallskip
    \textbf{Language:} Kotlin\\
    \textbf{Keywords:} Domain Specific Language, Aggregate Computing, Field Calculus, eXchange Calculus

    \divider

    \Needspace{10\baselineskip}
    \textcolor{accent}{\textbf{Project-Emerge}}\\
    \github{\href{https://github.com/Project-Emerge/Project-Emerge-system}{https://github.com/Project-Emerge/Project-Emerge-system}}\\
    \textit{06--09/2025}\\ \smallskip
    Open-source toolchain for multi-robot demonstrations based on aggregate computing.\\ \smallskip
    \textbf{Language:} Scala, Python, React\\
    \textbf{Keywords:} Aggregate Computing, Field Calculus, Swarm Robotics, Multi-Robot Systems, Web Application
    
    \divider

    \textcolor{accent}{\textbf{Rust Fields}}\\
    \github{\href{https://github.com/RustFields}{github.com/RustFields}}\\
    \textit{05--09/2023}\\ \smallskip
    Rust execution platform for ScaFi and field calculus on resource-constrained devices, developed across three Master's courses.\\ \smallskip
    \textbf{Language:} Rust, Scala\\
    \textbf{Keywords:} Aggregate Computing, Field Calculus, Rust, ScaFi, Pervasive Computing

    \divider

    \textcolor{accent}{\textbf{HarmoniKt}}\\
    \github{\href{https://github.com/HarmoniKt}{github.com/HarmoniKt}}\\
    \textit{05--08/2025}\\ \smallskip
    Microservices middleware providing standardised APIs for Boston Dynamics Spot and Mobile Industrial Robots (MiR),
    developed for the Ph.D. course ``Service Orchestration and Industrial IoT Platforms for Industry 4.0 and 5.0 Environments''.\\ \smallskip
    \textbf{Language:} Kotlin, Python\\
    \textbf{Keywords:} Microservices, Middleware, Robot Management

    \divider

    \cvgroupheading{Selected Academic and Personal Projects}

    \textcolor{accent}{\textbf{Equilessons}}\\
    \github{\href{https://github.com/angelacorte/equilessons}{https://github.com/angelacorte/equilessons}}\\
    \textit{2021/2023}\\ \smallskip
    Web application for managing riding-school lessons, started as a personal project and later
    extended for the Master's course ``Applicazioni e Servizi Web''.\\ \smallskip
    \smallskip

    \textbf{Language:} TypeScript\\
    \textbf{Keywords:} Web Application, Angular, MEAN, TypeScript

    \divider

    \textcolor{accent}{\textbf{Sette e Mezzo clone}}\\
    \github{\href{https://github.com/angelacorte/SetteMezzo-Clone}{github.com/angelacorte/SetteMezzo-Clone}}\\
    \textit{08--10/2022}\\ \smallskip
    Distributed CLI implementation of Sette e Mezzo for the Master's course ``Distributed Systems''.\\ \smallskip
    \textbf{Language:} TypeScript\\
    \textbf{Keywords:} CLI application, Distributed Systems, TypeScript

    \divider

    \textcolor{accent}{\textbf{Smart Charging Stations}}\\
    \github{\href{https://github.com/angelacorte/smart-charging-station-report}{github.com/angelacorte/smart-charging-station-report}}\\
    \textit{07--09/2023} \\ \smallskip
    EV charging-station management application for the Master's course ``Smart City e Tecnologie Mobili''.\\ \smallskip
    \textbf{Language:} Scala, TypeScript, Svelte, Node.js\\
    \textbf{Keywords:} Web Application, Smart City, Scala, TypeScript, Svelte, Akka Actors

    \divider

    \textcolor{accent}{\textbf{PPS-galaxy sim}}\\
    \github{\href{https://github.com/FilippoVissani/PPS-22-galaxy-sim}{github.com/FilippoVissani/PPS-22-galaxy-sim}}\\
    \textit{07--09/2023} \\ \smallskip
    Galaxy motion simulator developed for the Master's course ``Paradigmi di Programmazione e Sviluppo''.\\ \smallskip
    \textbf{Language:} Scala\\
    \textbf{Keywords:} Simulation, Scala, Akka Actors

\end{adjustwidth}

\end{document}
